\section{Baseline algorithms}
\label{sec:baselines}

\begin{lede}
The two published lifelong algorithms TOLL-PIBT is compared against, and the
space-time search both are built on. \Cref{sec:related} says what they are
for; this section says how they are implemented here.
\end{lede}

Both reuse the identical \code{assignment.TaskAssigner}
(\Cref{sec:assignment}), so only the low-level movement and
collision-avoidance layer differs. Neither uses candidate scoring, aisle
direction or reservations; both navigate with their own time-expanded,
reservation-table A$^\ast$, built from scratch because nothing else in this
codebase is collision-aware in that sense.

\subsection{Time-expanded \texorpdfstring{A$^\ast$}{A*}}
\label{sec:baselines-astar}

The shared primitive: cooperative A$^\ast$ over states $(\text{vertex}, t)$
rather than plain vertices, since avoiding other robots requires reasoning
about \emph{when} a cell is occupied, not just whether it is passable.
\begin{equation}
  f(v, t) = g(v, t) + h(v),
  \qquad
  h(v) = \Dg{\mathrm{goal}}(v),
\end{equation}
with $h$ the precomputed BFS distance map of \Cref{sec:graph-distances}. On an
unweighted grid where every move costs exactly 1, $h$ is the true shortest-path
distance and therefore admissible and consistent, so the search is optimal for
the single agent. Each transition is a move to a neighbour or a wait, each
costing $+1$, filtered through the reservation table --- which checks both a
vertex hold and, critically, a reverse-edge hold, catching swap conflicts the
same way \Cref{sec:pibt-conflicts} does.

Two modes:
\begin{itemize}
\item \code{require\_goal=True} (Token Passing): return only a path that
  reaches the goal \emph{and} can hold there for the rest of the horizon, the
  standard cooperative-A$^\ast$ goal condition, so a robot does not stop
  somewhere another robot will need later.
\item \code{require\_goal=False} (RHCR): the horizon is a short rolling window
  decoupled from goal distance, so reaching the goal within one window is the
  exception. Returns the best-effort path to whichever boundary state has the
  smallest $h$, ties broken by pop order.
\end{itemize}

\code{prioritized\_plan} plans robots one at a time, committing each path into
the shared table before planning the next --- the classical sequential MAPF
decomposition: no joint-optimality guarantee, always fast, always produces some
legal plan. Every robot's current cell is pre-held for one extra timestep so a
later, not-yet-planned robot cannot path straight through a robot that has not
moved yet.

\code{resolve\_residual\_conflicts} is a post-hoc safety net run after either
baseline finishes: it repeatedly finds any residual vertex or swap conflict
among the committed next positions and downgrades the lower-priority robot to
hold in place, iterating until stable --- bounded by $n+1$ iterations, since
each pass removes at least one conflict.

\subsection{Token Passing}
\label{sec:baselines-tp}

Each step, any robot whose committed path is stale (start or goal mismatch, or
exhausted) is replanned via \code{prioritized\_plan} with
\code{require\_goal=True}, against a table pre-seeded by holding every
still-valid path of settled robots. A robot whose search fails waits in place,
incrementing \code{tp\_path\_not\_found}.

Token Passing as published has \textbf{no priority-inheritance or backtracking
mechanism}: a blocked robot cannot ask the robot in its way to move. This is
the structural reason --- not an artefact of this implementation --- for the
corridor gridlock reported in \Cref{sec:results-baselines}: once a queue forms
in a single-file corridor with no mechanism to displace anyone, it never
resolves. The counters \code{tp\_path\_not\_found} and \code{tp\_forced\_holds}
are exposed so that a zero-throughput run can be confirmed as genuine gridlock
rather than a crash or a silent no-op.

\subsection{RHCR}
\label{sec:baselines-rhcr}

Every \param{baseline\_replan\_period} steps (default 5) all robots are jointly
replanned over a \param{baseline\_window}-length horizon (default 10) via
windowed CBS:

\begin{algorithm}[H]
\caption{Windowed CBS (\code{baselines/rhcr.py})}
\begin{algorithmic}[1]
\State \textbf{root}: each robot's unconstrained time-expanded A$^\ast$ path
  (\code{require\_goal=False}); $\mathrm{cost} = \sum_i |\mathrm{path}_i|$
\While{the open list is non-empty}
  \State pop the lowest-cost node
  \State $c \gets$ first conflict in the node's paths within the window
    \Comment{shared vertex, or a reversed-edge swap}
  \If{no conflict} \Return the node's paths \EndIf
  \State branch into two children, each adding one of the two constraints from $c$
  \State replan only the newly-constrained robot in each child; push if feasible
\EndWhile
\end{algorithmic}
\end{algorithm}

Branching follows \citet{sharon2015cbs} exactly: a vertex conflict $(i, j, v,
t)$ gives children $\{i \text{ forbidden at } (v,t)\}$ and $\{j \text{
forbidden at } (v,t)\}$; an edge conflict gives the two directed-edge
prohibitions. If node expansions exceed \param{node\_expansion\_cap} (500)
before a conflict-free node is found, the search falls back to
\code{prioritized\_plan} for that window and increments \code{cbs\_fallbacks}
--- standard, documented practice for bounding CBS's worst case, not a shortcut
unique to this implementation.

Between periodic replans, \code{\_repair\_stale\_paths} gives any robot with a
newly assigned task or an exhausted path an immediate single-agent repair,
rather than making it wait up to a full replan period.
